\documentclass[12pt]{article}
\usepackage{amsmath, amssymb, amsthm}
\usepackage{geometry}
\usepackage{hyperref}
\usepackage{booktabs}
\usepackage{graphicx}
\usepackage{listings}
\usepackage{xcolor}
\geometry{a4paper, margin=2.5cm}

\lstset{
    language=Python,
    basicstyle=\ttfamily\small,
    keywordstyle=\color{blue},
    commentstyle=\color{gray},
    stringstyle=\color{red},
    showstringspaces=false,
    frame=single,
    breaklines=true
}

\hypersetup{
    colorlinks=true,
    linkcolor=blue,
    urlcolor=blue,
    citecolor=blue
}

\title{The Hydrogen Atom from the Lagrangian Duality of the Fine-Structure Constant: \\
Independent Derivation via a Circular Matrix Product State}
\author{Massimiliano Blandino \\ 
\href{https://orcid.org/0009-0006-3252-4011}{ORCID: 0009-0006-3252-4011}}
\date{April 2026}

\begin{document}

\maketitle

\begin{abstract}
In previous works of this series we have demonstrated the duality between a geometric Lagrangian (oscillating circle of radius \(R = \pi\)) and the Polyakov Lagrangian for a bosonic string compactified on a circle of radius \(R = \pi\). The central result of this duality is:

\[
\alpha^{-1} = \ln(\lambda_{\max}) - \pi, \tag{1}
\]

where \(\lambda_{\max}\) is the dominant eigenvalue of a circular Matrix Product State (MPS) with bond dimension \(D = 45\).

In this independent paper we interpret the vacuum persistence amplitude of the compactified Polyakov string as:

\[
\Psi_{\text{vac}} = e^{-S_{\text{eff}}} = e^{-\alpha^{-1}} \approx e^{-137.036}, \tag{2}
\]

where \(S_{\text{eff}} = \ln\lambda_{\max} - \pi\) is the effective action.

This amplitude fixes the dimensionless coupling that governs the hydrogen atom. Using the universal constants \(\hbar, c, m_e\) (which set the lepton mass scale) we compute the physical properties of the hydrogen ground state. The dimensionless fine-structure constant \(\alpha\) is obtained from the MPS without any free parameters; the dimensional quantities (binding energy, Bohr radius, orbital velocity) follow from the electron mass scale. No additional free parameters are introduced.

We provide a complete, self-contained Python code, executable on Google Colab, that takes only the ontological parameters of the string (\(\pi,\ D=45,\ T=8,\ 24\) transverse modes) and outputs the hydrogen properties.

The hydrogen atom emerges as the ground state of a circular tensor network. The fine-structure constant is not a free parameter: it is the fundamental resonance of the vacuum. The construction is open, reproducible, and falsifiable.
\end{abstract}

\section{Introduction}

In previous chapters of this series we established the duality:

\[
\boxed{\alpha^{-1} = \ln(\lambda_{\max}) - \pi}, \tag{3}
\]

where \(\lambda_{\max}\) is the dominant eigenvalue of a circular MPS with bond dimension \(D = 45\). The duality originates from the equivalence between the geometric Lagrangian \(\mathcal{L}_{\text{geo}} = 4\dot{d}^2 + \frac{1}{R}d^2 + \frac{1}{4}\sqrt{R^2-d^2}\) (describing an oscillating circle of radius \(R = \pi\)) and the Polyakov Lagrangian for a bosonic string compactified on a circle of radius \(R = \pi\) \cite{polyakov1981, green2012, polchinski1998}.

The parameters that enter the duality are not free: they are derived from geometric identities and from the observed bound \(q \le 45\) on the continued fraction quotients of \(\alpha^{-1}\) over all historical CODATA values \cite{codata2022}. Specifically:
\begin{itemize}
    \item \(R = \pi\) (compactification radius);
    \item \(D = 45\) (bond dimension, maximum continued fraction quotient);
    \item \(T = 8\) (string tension, from \(\frac{T R^2}{2} = 4R^2\));
    \item \(N_{\text{trans}} = 24\) (transverse modes of the bosonic string, responsible for the Casimir effect);
    \item The curvature correction \(\mathcal{L}_{\text{curv}} = -1/(48\pi A(\pi))\) (where \(A(\pi)=4\pi^3+\pi^2+\pi\)).
\end{itemize}

In this independent paper we focus on the physical interpretation of the duality: the vacuum persistence amplitude of the compactified string, given by \(e^{-\alpha^{-1}}\), determines the scale of electromagnetic interactions. When combined with the lepton mass scale (set by \(m_e\)), this amplitude yields all the properties of the hydrogen atom.

\section{The Vacuum Persistence Amplitude and the String Wavefunction}

\subsection{From the MPS to the Vacuum Amplitude}

In the Polyakov-MPS duality, the effective action of the system is

\[
S_{\text{eff}} = \ln \lambda_{\max}. \tag{4}
\]

The closed string carries a topological phase \(\pi\) associated with the circumference of the circle. Subtracting this phase gives the effective action for the physical vacuum:

\[
S_{\text{phys}} = \ln \lambda_{\max} - \pi = \alpha^{-1}. \tag{5}
\]

The vacuum persistence amplitude (the probability that the vacuum remains in its ground state) is

\[
\Psi_{\text{vac}} = e^{-S_{\text{phys}}} = e^{-\alpha^{-1}}. \tag{6}
\]

This quantity is dimensionless and pure; it is not a spatial wavefunction. It represents the statistical weight of the entangled 45-node MPS configuration that gives rise to stable matter.

\subsection{Connection to the Hydrogen Atom}

The hydrogen atom is the simplest bound state of quantum electrodynamics. Its ground state properties are determined by the fine-structure constant \(\alpha\) and the electron mass \(m_e\). The MPS provides \(\alpha\) without any free parameters. The dimensional scale (eV, pm, km/s) is introduced by the universal constants \(\hbar, c, m_e\). Specifically:

\[
E_0 = -\frac{1}{2} m_e c^2 \alpha^2, \qquad
a_0 = \frac{\hbar}{m_e c \alpha}, \qquad
v = \alpha c. \tag{7}
\]

Thus, once \(\alpha\) is known, the hydrogen atom is completely determined relative to the lepton mass scale. The vacuum amplitude \(\Psi_{\text{vac}} = e^{-\alpha^{-1}}\) is not \(\psi_{1s}(0)\); rather, it is the weight that selects the vacuum configuration responsible for the Coulomb interaction.

\section{The Curvature Correction and the Number 24}

The curvature correction to the Lagrangian density is

\[
\mathcal{L}_{\text{curv}} = -\frac{1}{48\pi A(\pi)}. \tag{8}
\]

Its contribution to the action is obtained by integrating over the circle:

\[
S_{\text{curv}} = \int_0^{2\pi} \mathcal{L}_{\text{curv}} \, d\theta
= 2\pi \cdot \left(-\frac{1}{48\pi A(\pi)}\right)
= -\frac{1}{24A(\pi)}. \tag{9}
\]

Thus the factor \(48\pi\) in the Lagrangian density, after integration over the topological cycle, becomes the invariant \(24\) in the action. This \(24\) is the number of transverse modes of the bosonic string (Casimir effect) and also the third derivative of the polynomial \(A(x) = 4x^3 + x^2 + x\) evaluated at \(x = \pi\): \(A'''(\pi) = 24\).

\section{The Hydrogen Properties from the MPS Code}

\subsection{Algorithm Description}

Below is a complete Python code that implements the MPS contraction and computes the hydrogen properties. The code takes as input only the ontological parameters of the string (\(\pi, D=45, T=8, 24\) transverse modes) and outputs the hydrogen numbers. No experimental values of \(\alpha\) or any other physical constant (except \(\hbar, c, m_e, e\) for unit conversion) are used as fitting parameters. The code is executable on Google Colab.

\begin{lstlisting}[language=Python, caption=Hydrogen from the MPS (no free parameters)]
# -*- coding: utf-8 -*-
"""
Hydrogen atom from a circular Matrix Product State (MPS).

This script computes the fine-structure constant inverse alpha^{-1}
from the duality alpha^{-1} = ln(lambda_max) - pi, where lambda_max
is the dominant eigenvalue of the MPS transfer matrix constructed
from the geometric Lagrangian L_geo and the curvature correction L_curv.

The MPS parameters are:
    R = pi                (compactification radius)
    D = 45                (bond dimension, maximum continued fraction quotient)
    T = 8                 (string tension, from Polyakov duality)
    N_transverse = 24     (transverse bosonic string modes, Casimir effect)

No free parameters are used. The script outputs alpha^{-1} and the
corresponding hydrogen ground state properties (binding energy,
Bohr radius, orbital velocity, vacuum decay probability).

Author: Massimiliano Blandino
ORCID: https://orcid.org/0009-0006-3252-4011
Concept DOI: 10.5281/zenodo.19802606
"""

import numpy as np
from scipy.linalg import expm, eig

# ============================================================================
# 1. MPS PARAMETERS (ontological, not fitted)
# ============================================================================
R = np.pi                         # compactification radius
D = 45                            # bond dimension (max continued fraction quotient)
N_TRANSVERSE = 24                 # transverse bosonic string modes
T = 8                             # string tension (from Polyakov duality)
A_pi = 4*R**3 + R**2 + R          # polynomial 4π³ + π² + π

# Curvature correction (Casimir effect)
L_curv = -1.0 / (48.0 * np.pi * A_pi)

# Number of discretization steps (optimal from convergence analysis)
N_steps = 5000

# ============================================================================
# 2. GEOMETRIC LAGRANGIAN (oscillating circle, Chapter 3)
# ============================================================================
def L_geo(theta: float) -> float:
    """
    Geometric Lagrangian density for the oscillating circle.
    
    Args:
        theta: phase angle (0 to 2π)
    
    Returns:
        Lagrangian value at theta.
    """
    d = R * np.sin(theta)          # vertical displacement
    d_dot = R * np.cos(theta)      # displacement derivative
    chord = 2.0 * np.sqrt(max(R**2 - d**2, 0.0))
    # Kinetic + potential + surface term
    return 4.0 * d_dot**2 + (1.0 / R) * d**2 + 0.25 * chord

# ============================================================================
# 3. MPS TRANSFER MATRIX AND DOMINANT EIGENVALUE
# ============================================================================
def compute_lambda_max(N_steps: int) -> float:
    """
    Compute the dominant eigenvalue lambda_max of the circular MPS
    transfer matrix product.
    
    The continuous generator G(theta) = (L_geo(theta) + L_curv) * I_D
    is exponentiated to give the local transfer matrix.
    
    Args:
        N_steps: number of discretization steps along the circle.
    
    Returns:
        lambda_max: dominant eigenvalue of the product.
    """
    theta_vals = np.linspace(0.0, 2.0 * np.pi, N_steps, endpoint=False)
    d_theta = theta_vals[1] - theta_vals[0]
    M = np.eye(D, dtype=complex)      # total transfer matrix
    
    for theta in theta_vals:
        G_val = L_geo(theta) + L_curv
        G = G_val * np.eye(D)
        T_local = expm(G * d_theta)
        M = T_local @ M
    
    eigenvalues = eig(M, left=False, right=False)
    lambda_max = np.max(np.real(eigenvalues))
    return lambda_max

# ============================================================================
# 4. FINE-STRUCTURE CONSTANT FROM DUALITY
# ============================================================================
lambda_max = compute_lambda_max(N_steps)
ln_lambda_max = np.log(lambda_max)
alpha_inv = ln_lambda_max - np.pi
alpha = 1.0 / alpha_inv

# ============================================================================
# 5. HYDROGEN GROUND STATE PROPERTIES
# ============================================================================
# Physical constants (CODATA 2022, SI units)
c = 299792458.0                 # speed of light [m/s]
m_e = 9.1093837015e-31          # electron mass [kg]
hbar = 1.054571817e-34          # reduced Planck constant [J·s]
e_charge = 1.602176634e-19      # elementary charge [C]
eV = e_charge                   # 1 eV in Joules

# Binding energy (ground state, non-relativistic)
E0_J = -0.5 * m_e * c**2 * alpha**2
E0_eV = E0_J / eV

# Bohr radius
a0 = hbar / (m_e * c * alpha)
a0_pm = a0 * 1.0e12             # convert to picometers

# Orbital velocity (Bohr model)
v_orbit = alpha * c
v_kms = v_orbit / 1000.0        # convert to km/s

# Vacuum decay probability (amplitude of the string wavefunction)
P_decay = np.exp(-alpha_inv)

# ============================================================================
# 6. OUTPUT
# ============================================================================
print("=" * 60)
print("HYDROGEN FROM THE CIRCULAR MPS (NO FREE PARAMETERS)")
print("=" * 60)
print(f"\nMPS parameters:")
print(f"  Compactification radius R   = π = {np.pi:.6f}")
print(f"  Bond dimension D            = {D}")
print(f"  Transverse modes            = {N_TRANSVERSE}")
print(f"  String tension T            = {T}")
print(f"  Discretization steps N      = {N_steps}")
print(f"  ln(lambda_max)              = {ln_lambda_max:.12f}")
print(f"  pi                          = {np.pi:.12f}")
print(f"\nDuality: alpha^{-1} = ln(lambda_max) - pi")
print(f"  alpha^-1                    = {alpha_inv:.12f}")
print(f"  alpha                       = {alpha:.12f}")
print(f"\nHydrogen ground state properties:")
print(f"  Binding energy E0           = {E0_eV:.4f} eV")
print(f"  Bohr radius a0              = {a0_pm:.2f} pm")
print(f"  Orbital velocity v          = {v_kms:.2f} km/s")
print(f"  Vacuum decay probability    = {P_decay:.2e}")
print("=" * 60)
\end{lstlisting}

\subsection{Numerical Results}

Executing the code with \(N_{\text{steps}} = 5000\) yields:

\begin{verbatim}
============================================================
HYDROGEN FROM THE CIRCULAR MPS (NO FREE PARAMETERS)
============================================================

MPS parameters:
  Compactification radius R   = π = 3.141593
  Bond dimension D            = 45
  Transverse modes            = 24
  String tension T            = 8
  Discretization steps N      = 5000
  ln(lambda_max)              = 140.177591546953
  pi                          = 3.141592653590

Duality: alpha^-1 = ln(lambda_max) - pi
  alpha^-1                    = 137.035998893363
  alpha                       = 0.007297352579

Hydrogen ground state properties:
  Binding energy E0           = -13.6057 eV
  Bohr radius a0              = 52.92 pm
  Orbital velocity v          = 2187.69 km/s
  Vacuum decay probability    = 3.06e-60
============================================================
\end{verbatim}

These values coincide with experimental measurements within the numerical precision of the MPS contraction (dominated by float64 rounding).

\subsection{Interpretation of the Vacuum Decay Probability}

The vacuum decay probability \(P_{\text{decay}} = e^{-\alpha^{-1}} \approx 3 \times 10^{-60}\) is the probability per unit time (in natural units) that the entangled MPS configuration disintegrates back to the bare vacuum. Its extreme smallness explains why hydrogen atoms are effectively stable over cosmological timescales. The large value of \(\alpha^{-1} \approx 137\) (and hence the smallness of \(P_{\text{decay}}\)) is a direct consequence of the geometry of the string (\(R = \pi\)), its tension (\(T = 8\)), the 24 transverse Casimir modes, and the 45 nodes of the MPS.

\section{Discussion}

\subsection{Why 24? The Casimir Effect}

The number 24 appears in two distinct but related roles:
\begin{enumerate}
    \item As the third derivative of the polynomial \(A(x)=4x^3+x^2+x\) evaluated at \(x=\pi\): \(A'''(\pi)=24\). This corrects the action by \(-1/(24A(\pi))\).
    \item As the number of transverse modes of the bosonic string in 26 dimensions (after gauge fixing). This multiplies the microscopic polarizability to give the macroscopic dielectric constant.
\end{enumerate}
Both roles originate from the same quantum effect: the Casimir energy of a closed string. The integration of \(\mathcal{L}_{\text{curv}}\) over the circle explicitly links the factor \(48\pi\) in the Lagrangian density to the invariant \(24\) in the action (see Eq.~(9)).

\subsection{Why 45? The Bound on Continued Fraction Quotients}

All historical CODATA values of \(\alpha^{-1}\) produce continued fractions whose partial quotients do not exceed 45. This empirical bound is not arbitrary: it is the maximum bond dimension that sustains a stable resonance of the circular MPS. A quotient \(q > 45\) would break the spectral coherence and destroy the hydrogen atom.

\subsection{The Role of the Lepton Mass Scale}

The MPS determines the dimensionless fine-structure constant \(\alpha\) without free parameters. The dimensional properties of hydrogen (binding energy, Bohr radius, orbital velocity) require the electron mass \(m_e\) and the constants \(\hbar, c\). The electron mass is an external input that sets the scale of leptonic matter. The MPS does not predict \(m_e\); it predicts \(\alpha\). Once \(m_e\) is given, the entire architecture of the hydrogen atom is fixed relative to that mass scale.

\section{Conclusions}

We have presented an independent derivation of the hydrogen atom from the circular MPS that encodes the duality \(\alpha^{-1} = \ln(\lambda_{\max}) - \pi\). The key points are:

\begin{enumerate}
    \item The vacuum persistence amplitude \(\Psi_{\text{vac}} = e^{-\alpha^{-1}}\) is a dimensionless weight that selects the ground state configuration of the vacuum.
    \item The fine-structure constant \(\alpha\) is obtained from the MPS with bond dimension \(D=45\), string tension \(T=8\), and 24 transverse modes – all parameters are either geometric, topological, or derived from observation (the bound \(q \le 45\)).
    \item Given \(\alpha\) and the electron mass scale \(m_e\) (and the universal constants \(\hbar, c\)), the hydrogen ground state properties (binding energy \(-13.6057\ \text{eV}\), Bohr radius \(52.92\ \text{pm}\), orbital velocity \(2187.69\ \text{km/s}\)) emerge without any additional free parameters.
    \item The vacuum decay probability \(P_{\text{decay}} = e^{-137.036} \approx 3\times10^{-60}\) is a direct consequence of the large value of \(\alpha^{-1}\) and explains the extreme stability of hydrogen.
\end{enumerate}

The entire construction is open, reproducible, and falsifiable. The Python code provided can be run by any researcher on a standard laptop (via Google Colab) and reproduces the results without any hidden parameters.

The fine-structure constant is not a free parameter of nature. It is the resonance frequency of a string compactified on a circle of radius \(\pi\), projected onto a 3‑brane, and discretized into a 45‑node tensor network. The hydrogen atom is the ground state of this resonator – the most stable configuration of matter in the universe.

\section*{Acknowledgments}

During the preparation of this work, the author used two advanced language models for language revision, grammatical improvement, structural suggestions, and code documentation:

\begin{itemize}
    \item \textbf{DeepSeek-V3} (深度求索, DeepSeek Company, China)
    \item \textbf{Gemini Pro 3} (Google DeepMind, Alphabet Inc.)
\end{itemize}

Both models were used exclusively for non-scientific assistance. All scientific content – including the derivation of the duality \(\alpha^{-1} = \ln(\lambda_{\max}) - \pi\), the construction of the Lagrangians, the tensor network formulation, the falsification tests, and the hydrogen atom derivation – is the original work of the author. The author assumes full responsibility for the integrity, accuracy, and originality of this work.

\newpage

\bibliographystyle{plain}
\begin{thebibliography}{99}

\bibitem{polyakov1981}
Polyakov, A. M. (1981). Quantum geometry of bosonic strings.
\textit{Physics Letters B}, 103(3), 207-210.

\bibitem{green2012}
Green, M. B., Schwarz, J. H., \& Witten, E. (2012).
\textit{Superstring Theory: 25th Anniversary Edition}. Cambridge University Press.

\bibitem{polchinski1998}
Polchinski, J. (1998). \textit{String Theory, Vol. 1: An Introduction to the Bosonic String}.
Cambridge University Press.

\bibitem{verstraete2004}
Verstraete, F., \& Cirac, J. I. (2004). Matrix product states for quantum simulation.
\textit{Physical Review A}, 70(6), 062324.

\bibitem{schollwock2011}
Schollwöck, U. (2011). The density-matrix renormalization group in the age of matrix
product states. \textit{Annals of Physics}, 326(1), 96-192.

\bibitem{bohr1913}
Bohr, N. (1913). On the constitution of hydrogen atoms.
\textit{Philosophical Magazine}, 26(151), 1-25.

\bibitem{schrodinger1926}
Schrödinger, E. (1926). Quantization as an eigenvalue problem.
\textit{Annalen der Physik}, 384(4), 361-376.

\bibitem{codata2022}
CODATA Recommended Values of the Fundamental Physical Constants: 2022.
\href{https://physics.nist.gov/constants}{physics.nist.gov/constants}

\bibitem{alpha_series}
Blandino, M. (2026). The Lagrangian Duality of the Fine-Structure Constant.
Concept DOI: \href{https://doi.org/10.5281/zenodo.19802606}{10.5281/zenodo.19802606}.

\end{thebibliography}

\end{document}